In a fine-beam tube, electrons accelerated by a voltage of \UO\ make a thin glowing beam in a low-pressure gas. Helmholtz coils make a uniform magnetic field of \BO\ perpendicular to the beam; the beam forms a circle with a radius of \rO.

\begin{abcliste}
\abc Why is the beam a circle?

\abc Calculate the specific charge $e/m$ of the electron from these values.

\abc How does the circle change if the accelerating voltage is raised? Explain.
\end{abcliste}
